\documentclass[a4paper, 12pt]{article}

\usepackage[brazilian]{babel}
\usepackage{ae,aecompl}
\usepackage[utf8]{inputenc}

\usepackage{fullpage}  
\usepackage{indentfirst}      % Indenta o primeiro paragrafo das sections
\usepackage{graphicx}         % Utilizado para inserir imagens
\usepackage{hyperref}         % Cria link entre o sumario e as sections
\usepackage{fancyhdr}         % Permite alterar o header e footer
\usepackage[includefoot]{geometry}

% Seta o estilo padrao
\pagestyle{fancy}

% Configura links
\hypersetup{colorlinks=true,linkcolor=black,citecolor=black}

% Seta espacamento
\addtolength{\parskip}{2mm}
\addtolength{\footskip}{15mm}

% Coloca no pé a pagina atual
\begin{document}

% Inclua as variáveis
\newcommand{\titulo}{Documento de Projeto}
\newcommand{\subtitulo}{Produto: DIGI-DOC}
\newcommand{\autor}{Sérgio Oliveira Campos}
\newcommand{\data}{05 de Maio de 2008}
\newcommand{\cidade}{Brasília}

% Numero do contrato caso seja consultor
\newcommand{\ncontrato}{1234567890-32}


% Inclua header
\lhead {
	\setlength{\unitlength}{5mm} 
	\begin{picture}(0,0) 
		\includegraphics[scale=0.75]{img/header.pdf} 
	\end{picture} 
	\textsf{\vspace{1cm}}
}

\chead{}
\rhead{}
\renewcommand{\headrulewidth}{0pt}


% Inluia a capa
\newpage
% $Id$
% ---------------------------------------------------------------------------
%
%  This is part of the relatorio-modelo.
%  Copyright (C) 2008 Interlegis
%  See the file relatorio.tex for copying conditions.
%

\textsf{\vspace{6cm}}
\begin{center}
  \noindent
  \huge{
    \textbf{\titulo}
  } \\
  \Large{
    \textbf{\subtitulo}
  }\\
  \large{
    \textbf{\subsubtitulo}
  }

  \vspace{9cm}

  \large{
    \textbf{\autor}\\
    \textsf{Contrato N$^{\circ}$: \numcontrato}\\
  }
\end{center}
\cfoot{}                        % retira o rodape

%
% Local variables:
%   mode: flyspell
%   TeX-master: "relatorio.tex"
% End:
%


% Seta o footer das paginas iniciais para numeros Romanos
\clearpage
\cfoot{\thepage}
\setcounter{page}{1}
\pagenumbering{Roman}

% Sumário
\newpage
\tableofcontents

% Lista de figuras
\newpage
\listoffigures

% Lista de tabelas
\newpage
\listoftables

% Seta o footer das paginas para numeros convencionais
\clearpage
\setcounter{page}{1}
\pagenumbering{arabic}

% PARA INSERIR NOVAS PAGINAS BASTA ADICIONAR AQUI:
% Conteúdo
\include{capit/1}
\include{capit/2}

\end{document}
